% TOP_PROBLEM: {"id": 30006733, "title": "The Haldane Conjecture on Spectral Gaps of Antiferromagnetic Quantum Spin Chains", "category_id": 16, "category": {"id": 16, "name": "physics", "display_name": "Mathematical Physics", "description": "Problems at the intersection of mathematics and physics.", "slug": "physics", "order_index": 16, "created_at": "2026-07-31T15:26:25.671Z"}, "status": "open"}
% FIRST_POSTED: 2026-09-27
% !TeX program = lualatex
% Compile twice: lualatex 30006733.tex
% All references and prose are embedded; no BibTeX database or external inputs.
\documentclass[11pt,a4paper]{article}
\usepackage[margin=24mm,headheight=14pt]{geometry}
\usepackage{fontspec}
\setmainfont{Latin Modern Roman}
\setsansfont{Latin Modern Sans}
\setmonofont{Latin Modern Mono}
\usepackage{amsmath,amssymb,mathtools}
\usepackage{unicode-math}
\setmathfont{Latin Modern Math}
\usepackage{microtype}
\usepackage{enumitem,xurl,needspace}
\usepackage[dvipsnames]{xcolor}
\usepackage{hyperref}
\hypersetup{unicode=true,colorlinks=true,linkcolor=MidnightBlue,urlcolor=MidnightBlue,
 pdftitle={Research attempt: Problem 30006733},pdfauthor={Research notebook},bookmarksnumbered=true}
\usepackage{bookmark}
\usepackage{tocloft}
\setlength{\cftsecnumwidth}{3em}
\cftsetpnumwidth{2.5em}
\cftsetrmarg{4em}
\usepackage{titlesec}
\titleformat{\section}{\Large\bfseries\raggedright}{\thesection}{0.7em}{}
\usepackage{fancyhdr}
\pagestyle{fancy}\fancyhf{}
\fancyhead[L]{\small\sffamily Open Problems Math | Research notebook}
\fancyhead[R]{\small Problem 30006733 | 27 September 2026}
\fancyfoot[C]{\thepage}
\setlength{\parindent}{0pt}
\setlength{\parskip}{3pt plus 1pt}
\setlength{\emergencystretch}{3em}
\setcounter{tocdepth}{1}
\setcounter{secnumdepth}{1}
\setlist[itemize]{leftmargin=3em,itemsep=5pt,topsep=4pt,parsep=0pt}
\allowdisplaybreaks
\newcommand{\N}{\mathbb{N}}
\newcommand{\Z}{\mathbb{Z}}
\newcommand{\Q}{\mathbb{Q}}
\newcommand{\R}{\mathbb{R}}
\newcommand{\C}{\mathbb{C}}
\newcommand{\F}{\mathbb{F}}
\newcommand{\A}{\mathbb{A}}
\newcommand{\PP}{\mathbb P}
\newcommand{\E}{\mathbb{E}}
\newcommand{\eps}{\varepsilon}
\newcommand{\dd}{\,\mathrm d}
\newcommand{\sref}[2]{\hyperlink{source:#1:#2}{[S#2]}}
\newcommand{\eref}[2]{\hyperlink{extra:#1:#2}{[E#2]}}
% Turn the next switch off for a shorter reading copy. The quotations preserve
% every exactTarget field, including qualifications omitted from a short summary.
\newif\ifshowcatalogue
\showcataloguetrue
\newcommand{\cataloguescope}[1]{\ifshowcatalogue
 \paragraph{Exact scope in the supplied catalogue.}
 \begin{quote}\small #1\end{quote}\fi}
\usepackage{etoolbox}
\AtBeginEnvironment{itemize}{\small}
\numberwithin{equation}{section}
% Standalone export. Original research content preserved.
\begin{document}
\setcounter{section}{0}
% ================================================================
% problemId: 30006733
% Canonical problem ID: 30006733
% exactTarget SHA-256: e4880dcda7df179eeffebbb8632036e6c4f0baf24444815395db6efff4049ae9
\clearpage
\section*{\texorpdfstring{The Haldane Conjecture on Spectral Gaps of Antiferromagnetic Quantum Spin Chains}{The Haldane Conjecture on Spectral Gaps of Antiferromagnetic Quantum Spin Chains}}\label{30006733}
\subsection{Definitions and mathematical statement}
A spin-$S$ site has Hilbert space $\C^{2S+1}$ and angular-momentum operators $S^x,S^y,S^z$. In the usual even-length periodic approximation, the antiferromagnetic chain is
\[
 H_L=J\sum_{j=1}^{L}\mathbf S_j\cdot\mathbf S_{j+1},\qquad J>0,\quad\mathbf S_{L+1}=\mathbf S_1.
\]
Let $\Delta_L$ be the energy difference between the ground energy and the next distinct energy. The proposed thermodynamic behavior is a positive gap for integer $S\ge1$, and gaplessness for $S\in\tfrac12+\Z_{\ge0}$. Periodic or infinite-chain conventions avoid confusing boundary edge modes with a bulk gap. The precise infinite-volume spectral statement is the one in the source.
\par\smallskip\textit{Formulation and concept references:} \sref{30006733}{1}.
\cataloguescope{Prove that the one-dimensional isotropic antiferromagnetic Heisenberg chain has a positive spectral gap for every integer spin S >= 1 and is gapless for every half-integer spin.}
\subsection{Short English statement}
Does the isotropic antiferromagnetic Heisenberg chain have a bulk energy gap exactly for integer, rather than half-integer, spin?
% BEGIN RESEARCH ADDITION 111
\subsection{Research attempt}

\paragraph{Current frontier and formulation audit.}
The full accessible manuscript corresponding to the original source is
arXiv:2409.09685v2. Its introduction explicitly separates the integer-spin
Heisenberg conjecture from the frustration-free setting of its theorems.
Its inverse-square finite-size criterion assumes frustration-freeness;
Remark~1.5 also warns about periodic versus open boundaries. Those hypotheses
cannot be transferred to the present Hamiltonian. \sref{30006733}{1}\eref{30006733}{1}

The half-integer obstruction is of Lieb--Schultz--Mattis type. A precise
infinite-chain theorem rules out a translation-invariant locally unique
gapped ground state under the stated symmetry assumptions; it is not, for
an arbitrary spin Hamiltonian, a theorem excluding symmetry-breaking
gapped phases. Tasaki's review carefully distinguishes these alternatives.
The source incorporated by the catalogue does not further specify which
infinite-volume ground-state representation its half-integer sentence uses.
The derivation below therefore proves the explicit even-periodic
finite-size assertion, and does not silently identify it with every possible
infinite-volume assertion. \eref{30006733}{2}

\paragraph{Attack route.}
Three routes were tested: an explicit twist trial state for half-integer
spin; operator comparison with the spin-one AKLT parent Hamiltonian; and a
continuous interpolation from that parent Hamiltonian. The first needs an
orthogonality argument, not just a low trial energy. The second needs control
of the actual ground energy and ground projection. The third needs a
size-independent gap estimate all along the path, rather than continuity
at each fixed chain length.

\paragraph{Attempt: a complete finite-volume calculation.}
We take even $L\geq4$ and keep $J$ explicit. First we justify the ground-state
input instead of assuming it. In the product basis of $S_j^z$ eigenvectors,
restrict to total magnetization zero. Multiplication by
\[
 (-1)^{\sum_{j\ \mathrm{odd}}(S-m_j)}
\]
changes every nonzero off-diagonal exchange matrix element to a negative
number. The configuration graph in this sector is connected: successive
nearest-neighbor transfers of one unit of magnetization redistribute any
configuration with the same total into any other, never exceeding the
capacities $-S\leq m_j\leq S$. Equivalently, this is the connected configuration
graph of a fixed number of indistinguishable units on a connected cycle,
with capacity $2S$ at each site. Perron--Frobenius, applied to a sufficiently
large scalar minus this matrix, gives a unique sector ground vector with
strictly positive transformed coefficients.

The product of singlets on $(1,2),(3,4),\ldots,(L-1,L)$ has nonnegative
transformed coefficients and nonzero overlap with this vector. Since total
spin commutes with $H_L$, sector uniqueness makes the ground vector an
eigenvector of total spin squared. Its overlap with a singlet makes that
spin zero. Every total-spin multiplet on an even chain has integral spin
and contains a zero-magnetization vector. Consequently this sector ground
vector, denoted $\Omega$, is the unique global ground state and is a
singlet. Translation $T$ therefore satisfies $T\Omega=e^{i\vartheta}\Omega$.
These are the standard sign/positivity ingredients of the finite-chain
argument; the proof above states the needed version explicitly.

Put $q=2\pi/L$ and
\[
 U=\exp\left(iq\sum_{j=1}^L jS_j^z\right).
\]
With the convention $TS_j^zT^{-1}=S_{j+1}^z$,
\[
 TUT^{-1}=U\exp(-iqS_{\mathrm{tot}}^z)\exp(2\pi iS_1^z).
\]
For half-integer $S$, the last factor is $-I$. Thus both $U\Omega$ and
$U^{-1}\Omega$ have translation eigenvalue $-e^{i\vartheta}$ and are
orthogonal to $\Omega$. This also covers the boundary bond, since rotation
of a spin operator through $2\pi$ is the identity even when the unitary on
state vectors changes sign.

Write $A_j=S_j^xS_{j+1}^x+S_j^yS_{j+1}^y$. Averaging the two twists cancels
the sine/current terms exactly and gives
\[
 \frac12\sum_{\epsilon=\pm1}
 \langle U^\epsilon\Omega,(H_L-E_0)U^\epsilon\Omega\rangle
 =J(\cos q-1)\sum_j\langle\Omega,A_j\Omega\rangle.
\]
Since $\|A_j\|\leq2S^2$, at least one normalized orthogonal trial vector has
energy excess at most
\[
 2JS^2L(1-\cos(2\pi/L))\leq\frac{4\pi^2JS^2}{L}.
\]
The variational principle proves
\begin{equation}\label{eq:30006733:twist}
 0<\Delta_L\leq\frac{4\pi^2JS^2}{L}
 \qquad(S\in\tfrac12+\Z_{\geq0},\ L\geq4\text{ even}).
\end{equation}
The strict left inequality is finite-dimensional uniqueness, not a uniform
lower bound. Equation~\eqref{eq:30006733:twist} settles the gap-closing conclusion
in the stated periodic approximation for every fixed half-integer $S$.
It is a rederivation of an established mechanism, not a new Haldane proof.
\eref{30006733}{2}

\paragraph{Attempt: why a positive AKLT comparison fails.}
For spin one, let $Q_j=\mathbf S_j\cdot\mathbf S_{j+1}$ and let $P_{a,j}$
project onto total bond spin $a=0,1,2$. The eigenvalues of $Q_j$ are
$-2,-1,1$, respectively. Direct polynomial interpolation gives
\[
 P_{2,j}=\frac{Q_j^2+3Q_j+2I}{6},\qquad
 Q_j+2I=P_{1,j}+3P_{2,j}.
\]
Consequently, with $J=1$ and $H_A=\sum_jP_{2,j}$,
\begin{equation}\label{eq:30006733:comparison}
 H_L+2LI=3H_A+\sum_jP_{1,j}\ \geq\ 3H_A.
\end{equation}
The spin-one AKLT parent is a standard gapped frustration-free example
within the source's framework. \eref{30006733}{1}
Nevertheless \eqref{eq:30006733:comparison} does not compare spectral gaps:
it compares unshifted energies relative to different ground spaces.
Already $A=\operatorname{diag}(0,1)$ and
$B=\operatorname{diag}(1,1+\delta)$, $0<\delta<1$, satisfy $B\geq A$,
while their gaps are $1$ and $\delta$.

In fact the shifted Heisenberg Hamiltonian is not frustration-free.
The kernel of $Q_j+2I$ is the one-dimensional bond singlet. A vector that is
a singlet on sites $1,2$ factors as that singlet tensor the remaining sites.
Its two-site density on $2,3$ is therefore $I_{\mathcal H_2}/(2S+1)\otimes\rho_3$, which
is not a pure singlet. Adjacent bond kernels have no common nonzero vector.
Thus the positive extensive ground energy in
\eqref{eq:30006733:comparison} is precisely what prevents the direct use of a
frustration-free gap criterion.

One can interpolate using
\[
 H_L(t)=3H_A+t\sum_jP_{1,j},\qquad0\leq t\leq1.
\]
At $t=1$ this is $H_L+2LI$. At fixed $L$, eigenvalue continuity is elementary,
but $\|H_L(t)-H_L(s)\|\leq L|t-s|$. Weyl's inequality therefore loses a
factor $L$ when controlling the gap. A small interval of length $O(L^{-1})$
near $t=0$ is not a uniform path to $t=1$.

Here is a precise sufficient replacement. Let $P_L(t)$ be the actual
ground projection and $E_L(t)$ the actual ground energy. If one could prove,
with $c>0$ independent of even $L$ and $t$, that
\begin{equation}\label{eq:30006733:path}
 \langle v,(H_L(t)-E_L(t))v\rangle\geq c\|v\|^2
 \quad\text{for }v\perp\operatorname{ran}P_L(t),
\end{equation}
then the spin-one periodic gap would follow. The obstacle is not the
existence of the interpolation but the uniform estimate involving its
moving ground space. This route does not even address all integer spins
without further work.

\paragraph{Stress test and reproducibility.}
The companion script constructs the local spin matrices and periodic
Hamiltonians for selected small $L$ and $S$. Its eigenvalues are floating-point
numerical evidence, not certified thermodynamic bounds. The identities for
$P_2$ and the two-by-two gap-comparison counterexample are exact. No finite
list of positive gaps is used to infer a positive limit.

The twist proof uses uniqueness to ensure that an orthogonal trial state
lies above the ground energy. Without that input, orthogonality to a chosen
ground vector need not exclude the rest of the ground space, especially
because the catalogue defines the next \emph{distinct} energy. For integer
spin, $\exp(2\pi iS_1^z)=I$, so the same twist has no forced orthogonality.
Finally, small periodic splittings alone are not promoted to a statement
about the excitation gap above every pure infinite-chain state.
\eref{30006733}{2}

\paragraph{Outcome.}
\textbf{Outcome: Exploratory approach with unresolved gap.}

The periodic half-integer gap-closing estimate is proved above with explicit
constants and its ground-state hypothesis supplied. For integer spin, the
attempt identifies and disproves a tempting monotonicity argument and
isolates the missing size-independent moving-ground-space estimate.
Neither that estimate nor the full integer-spin claim has been established.
There is no claim of novelty for the half-integer branch or the AKLT facts.
% END RESEARCH ADDITION 111
\subsection{Sources}
\begin{itemize}\raggedright
\item[\textnormal{[S1]}]\hypertarget{source:30006733:1}{} Marius Lemm and Angelo Lucia, On the critical finite-size gap scaling for frustration-free Hamiltonians, Review of Mathematical Physics 37 (2025), introduction.
\url{https://re.public.polimi.it/retrieve/2332a1fb-6d7b-4ef7-a66f-1cbb552b76f8/lemm-lucia-2025-on-the-critical-finite-size-gap-scaling-for-frustration-free-hamiltonians.pdf}.
{\small\textit{Catalogue formulation source.}}
% BEGIN NEW REFERENCES 111
\item[\textnormal{[E1]}]\hypertarget{extra:30006733:1}{} Marius Lemm and Angelo Lucia,
\emph{On the critical finite-size gap scaling for frustration-free Hamiltonians},
arXiv:2409.09685v2 (2025), introduction, Theorem 1.4 and Remark 1.5;
accessible full manuscript of [S1].
\url{https://arxiv.org/html/2409.09685v2}.
\item[\textnormal{[E2]}]\hypertarget{extra:30006733:2}{} Hal Tasaki,
\emph{The Lieb--Schultz--Mattis Theorem: A Topological Point of View},
arXiv:2202.06243v2 (2022), introduction, Sections 2--3, especially the
finite/infinite-volume distinction and the twist-operator argument.
\url{https://arxiv.org/html/2202.06243v2}.
% END NEW REFERENCES 111
\end{itemize}

\end{document}
